\documentclass{ximera}

\input{../../preamble.tex}

\title{Roots and Radicals}

\begin{document}

\begin{abstract}
aspects
\end{abstract}
\maketitle





Polynomial and rational functions only include terms with integer powers. 


On the other hand, power functions use any real number power.  The noninteger powers belong to a category called \textbf{roots and radicals}.

We saw this with our Core Root and Radical functions.



\begin{definition} \textbf{\textcolor{green!50!black}{Roots}}

$x^{\tfrac{1}{n}}$ is called \textbf{the nth root of x}, where $n$ is a natural number.




\textbf{odd roots}

When $n = 2k+1$ is odd, $x^{\tfrac{1}{n}} = x^{\tfrac{1}{2k+1}}$ is defined to be the unique real number, $r$, such that $r^n = x$. The domain is all real numbers.




\textbf{even roots}

When $n=2k$ is even, $x^{\tfrac{1}{n}} = x^{\tfrac{1}{2k}}$ is defined to be the unique nonnegative real number, $r$, such that $r^n = x$. The domain is all nonnegative real numbers.


\end{definition}



\begin{definition} \textbf{\textcolor{green!50!black}{Radicals}}

An alternate point-of-view uses radical notation rather than exponents.

\[   \sqrt[n]{x} =  x^{\tfrac{1}{n}}     \]

where $\sqrt{ }$ is called the \textbf{radical} symbol.


$\sqrt[n]{x}$ is called the \textbf{nth root of x}.

\end{definition}




\subsection*{Algebra Rules}

The nth-root is an exponent, so it follows all of the exponent rules.


\begin{itemize}
\item $a^n \cdot a^m = a^{n+m}$

\item $\frac{a^n}{a^m} = a^{n-m}$

\item $(a^n)^m = a^{n \cdot m}$

\item $a^n \cdot b^n = (a \cdot b)^n$

\item $\frac{a^n}{b^n} = \left(\frac{a}{b}\right)^n$


\end{itemize}



Similar rules for radicals.

\begin{itemize}

\item $\sqrt[n]{a} \cdot \sqrt[n]{b} = \sqrt[n]{a \cdot b}$

\item $\frac{\sqrt[n]{a}}{\sqrt[n]{b}} = \sqrt[n]{\frac{a}{b}}$

\item $\sqrt[m]{\sqrt[n]{a}} = \sqrt[nm]{a}$

\end{itemize}





















\subsection*{Core Radical Functions}



Negative numbers raised to odd powers result in negative numbers, which means every real number has an odd root.  The domain of odd roots or radicals is all real numbers.

\begin{example}


The graph of $y = R(t) = t^{\tfrac{1}{3}} = \sqrt[3]{t}$


\begin{image}
\begin{tikzpicture} 
  \begin{axis}[
            domain=-10:10, ymax=10, xmax=10, ymin=-10, xmin=-10,
            axis lines =center, xlabel=$t$, ylabel=$y$,
            ytick={-10,-8,-6,-4,-2,2,4,6,8,10},
            xtick={-10,-8,-6,-4,-2,2,4,6,8,10},
            ticklabel style={font=\scriptsize},
            every axis y label/.style={at=(current axis.above origin),anchor=south},
            every axis x label/.style={at=(current axis.right of origin),anchor=west},
            axis on top
          ]
          
          \addplot [line width=2, penColor, smooth, samples=200, domain=(0:9),->] {x^(0.333333)};
          \addplot [line width=2, penColor, smooth, samples=200, domain=(-9:0),<-] {-(-x)^(0.333333)};



           

  \end{axis}
\end{tikzpicture}
\end{image}

\end{example}





Negative numbers raised to even powers result in positive numbers, just like positive numbers raised to even power.  This makes going backwards a problem.  



Should $\sqrt{4} = -2$ or $\sqrt{4} = 2$?  When squared, both $-2$ and $2$ give $4$.



We are forced to choose between two even roots.  We pick positive numbers . The domain and range of even roots are both $[0, \infty)$.

\begin{example}


The graph of $y = R(t) = t^{\tfrac{1}{2}} = \sqrt{t} $


\begin{image}
\begin{tikzpicture} 
  \begin{axis}[
            domain=-10:10, ymax=10, xmax=10, ymin=-10, xmin=-10,
            axis lines =center, xlabel=$t$, ylabel=$y$,
            ytick={-10,-8,-6,-4,-2,2,4,6,8,10},
            xtick={-10,-8,-6,-4,-2,2,4,6,8,10},
            ticklabel style={font=\scriptsize},
            every axis y label/.style={at=(current axis.above origin),anchor=south},
            every axis x label/.style={at=(current axis.right of origin),anchor=west},
            axis on top
          ]
          
          \addplot [line width=2, penColor, smooth, samples=200, domain=(0:9),->] {x^(0.5)};
          \addplot[color=penColor,fill=penColor,only marks,mark=*] coordinates{(0,0)};




           

  \end{axis}
\end{tikzpicture}
\end{image}

\end{example}






There are no horizontal asymptotes on these graphs.  The graph keeps moving up or down unbounded.  The root and radical functions are unbounded. Their values tend to $\pm \infty$ as the domain moves out towards $\pm \infty$.
















\subsection*{Root/Radical Functions}

The Elementary Functions include root and radical functions.

These are constructed from the core functions by composing them with linear functions.





\begin{definition} \textbf{\textcolor{green!50!black}{Root and Radical Functions}}


Let $R(k) = \sqrt[n]{k} =  k^{\tfrac{1}{n}} $ be a Core root function.

Let $L_{in}(t) = B \, t + C$ be a linear function.

Let $L_{out}(y) = A \, y + D$ be a linear function.


Root and Radical functions are constructed by composing these component functions.


\[
(L_{out} \circ R \circ L_{in})(x) = A \sqrt[n]{B \,x + C} + D
\]




\end{definition}






























\subsection*{Even Roots}

This course is the study of the real numbers. As a result we don't take even roots of negative numbers, because a negative number raised to an even power always results in a positive number.



\[   EvenRoot(r) = \sqrt[2n]{r} = r^{\tfrac{1}{2n}}          \]

The square root function is the most popular core root function


\[   \sqrt{x} = x^{\tfrac{1}{2}}          \]



The domain of the square root function (and all even roots) is $[0, \infty)$, the \textbf{nonnegative} numbers.  






The graph of $y = SQR(x) = \sqrt{x}$


\begin{image}
\begin{tikzpicture} 
  \begin{axis}[
            domain=-10:10, ymax=10, xmax=10, ymin=-10, xmin=-10,
            axis lines =center, xlabel=$x$, ylabel=$y$,
            ytick={-10,-8,-6,-4,-2,2,4,6,8,10},
            xtick={-10,-8,-6,-4,-2,2,4,6,8,10},
            ticklabel style={font=\scriptsize},
            every axis y label/.style={at=(current axis.above origin),anchor=south},
            every axis x label/.style={at=(current axis.right of origin),anchor=west},
            axis on top
          ]
          
            %\addplot [line width=2, penColor, smooth, domain=(-9:0),<->] {(x-1)/((x+3)*(x-4))};
            \addplot [line width=2, penColor, smooth, samples=200,domain=(0:9),->] {sqrt(x)};
   
      \addplot[color=penColor,fill=penColor,only marks,mark=*] coordinates{(0,0)};

           

  \end{axis}
\end{tikzpicture}
\end{image}



The square root function begins where the inside of the radical equals $0$.  It then moves in the direction that keeps the inside positive.


All even root/radical functions can be viewed as the core square root function composed with linear function. One linear function affects the inside of the radical sign and the other affects the outside.


\begin{idea} \textbf{\textcolor{blue!55!black}{Square Root Functions}} 


Let $SQR(x) = \sqrt{k}$ be the core square root function.

Let $L_{in}(t) = B \, t + C$ be a linear function.

Let $L_{out}(y) = A \, y + D$ be a linear function.


In general, square root (radical) functions are constructed by composing these component functions.


\[
(L_{out} \circ SQR \circ L_{in})(x) = A \sqrt{B \,x + C} + D
\]


\end{idea}







\begin{example} Square Root Function

The graph of $y = f(v) = \frac{1}{3} \sqrt{5-3v}-4$


\begin{image}
\begin{tikzpicture} 
  \begin{axis}[
            domain=-10:10, ymax=10, xmax=10, ymin=-10, xmin=-10,
            axis lines =center, xlabel=$v$, ylabel=$y$, 
            ytick={-10,-8,-6,-4,-2,2,4,6,8,10},
            xtick={-10,-8,-6,-4,-2,2,4,6,8,10},
            ticklabel style={font=\scriptsize},
            every axis y label/.style={at=(current axis.above origin),anchor=south},
            every axis x label/.style={at=(current axis.right of origin),anchor=west},
            axis on top
          ]
          
            %\addplot [line width=2, penColor, smooth, domain=(-9:0),<->] {(x-1)/((x+3)*(x-4))};
            \addplot [line width=2, penColor, smooth, samples=200,domain=(-9:2.67),<-] {0.33*sqrt(5-3*x)-4};
   
      \addplot[color=penColor,fill=penColor,only marks,mark=*] coordinates{(1.67,-4)};

           

  \end{axis}
\end{tikzpicture}
\end{image}

\textbf{Domain:} The domain of $f$ is the interval that makes the inside of the radical nonnegative.  The inside is $5-3v$, which is a linear function with a negative leading coefficient. The negative leading coefficient tells us that $5-3v > 0$ on the negative side of its zero, which is $\frac{5}{3}$.

\[ \left(-\infty, \frac{5}{3} \right] \]

\textbf{Zeros:}


\begin{align*}
\frac{1}{3} \sqrt{5-3v}-4 & = 0   \\
\sqrt{5-3v}               & = 12  \\
5-3v                      & = 36  \\
-31 &= 3v   \\
-\frac{31}{3}   &= v
\end{align*}


\textbf{Continuity:} Root or Radical functions are continuous.


\textbf{End-Behavior:} Root or Radical functions have end-behavior in only one direction. The domain tells us which direction.  

The leading coefficient tells us the end-behaviore. Here the leading coefficient of $f$ is $\frac{1}{3}$, which is positive.


\[  \lim\limits_{v \to -\infty} f(v) = \infty   \]







\textbf{Behavior:}  We'll use the Chain Rule.

$f$ is the composition of

\begin{itemize}
\item $SQR(x) = \sqrt{x}$ 

\item $L_{in}(t) = 5 - 3 \, t$, a decreasing linear function

\item $L_{out}(y) = \frac{1}{3} \, y + 4$, an increasing linear function
\end{itemize}


\[  f(v) =  \frac{1}{3} \sqrt{5-3v}-4 = (L_{out} \circ SQR \circ L_{in})(v) \]

\[ inc \circ inc \circ dec = dec  \]


$f$ is a decreasing square root function, which agrees with the graph.



\textbf{Global Maximum and Minimum:} 


Since $\lim\limits_{v \to -\infty} f(v) = \infty$, we know $f$ has no global maximum.

Since $f$ is always decreasing, we know that $f\left( \frac{5}{3} \right) = -4$ is the global minimum.


\textbf{Local Maximums and Minimums:} 

Global extrema are automatically local extrema.  So, $f\left( \frac{5}{3} \right) = -4$ is a local minimum.

Since $f$ is always decreasing, there are no other local extrema.


\textbf{Range:} 

\begin{itemize}
\item $f$ is continuous
\item $\lim\limits_{v \to -\infty} f(v) = \infty$
\item $-4$ is the global minimum
\end{itemize}

All of this gives the range  $\left[ -4, \infty \right)$



\end{example}















\begin{example} Even Root

The graph of $y = T(k) = \sqrt{-k+5}$


\begin{image}
\begin{tikzpicture} 
  \begin{axis}[
            domain=-10:10, ymax=10, xmax=10, ymin=-10, xmin=-10,
            axis lines =center, xlabel=$k$, ylabel=$y$, 
            ytick={-10,-8,-6,-4,-2,2,4,6,8,10},
            xtick={-10,-8,-6,-4,-2,2,4,6,8,10},
            ticklabel style={font=\scriptsize},
            every axis y label/.style={at=(current axis.above origin),anchor=south},
            every axis x label/.style={at=(current axis.right of origin),anchor=west},
            axis on top
          ]
          
            %\addplot [line width=2, penColor, smooth, domain=(-9:0),<->] {(x-1)/((x+3)*(x-4))};
            \addplot [line width=2, penColor, smooth, samples=200,domain=(-9:5),<-] {sqrt(-x+5)};
   
      \addplot[color=penColor,fill=penColor,only marks,mark=*] coordinates{(5,0)};

           

  \end{axis}
\end{tikzpicture}
\end{image}


\textbf{Domain:} 
Here the inside $-k+5=0$ equals $0$ when $k=\answer{5}$.  That is the start of the domain.  Then the inside is positive $-k+5>0$, when $k$ \wordChoice{\choice[correct]{$<$} \choice {$>$}}  $5$, which means the domain is $(-\infty,5]$.


\textbf{Zeros:}  Root and Radical functions can only have at most one zero, which must be $5$.


\textbf{Continuity:} Root and Radical functions are continuous. 


\textbf{End-Behavior:} Root or Radical functions have end-behavior in only one direction. The domain tells us which direction.  

The leading coefficient tells us the end-behaviore. Here the leading coefficient of $T$ is $1$, which is positive.


\[  \lim\limits_{k \to -\infty} T(k) = \infty   \]







\textbf{Behavior:}  We'll use the Chain Rule.

$T$ is the composition of

\begin{itemize}
\item $SQR(x) = \sqrt{x}$, an increasing core function
\item $L_{in}(t) = -k + 5$, a decreasing linear function
\end{itemize}


\[  T(k) =  \sqrt{-k+5} = (SQR \circ L_{in})(k) \]

\[ inc \circ dec = dec  \]


$T$ is a decreasing square root function, which agrees with the graph.





\textbf{Global Maximum and Minimum:} 


Since $\lim\limits_{k \to -\infty} T(k) = \infty$, we know $T$ has no global maximum.

Since $T$ is always decreasing, we know that $T(5) = 0$ is the global minimum.


\textbf{Local Maximums and Minimums:} 

Global extrema are automatically local extrema.  So, $T(5) = 0$ is a local minimum.

Since $T$ is always decreasing, there are no other local extrema.


\textbf{Range:} 

\begin{itemize}
\item $T$ is continuous
\item $\lim\limits_{k \to -\infty} T(k) = \infty$
\item $0$ is the global minimum
\end{itemize}

All of this gives the range  $\left[ 0, \infty \right)$.






\end{example}









\begin{summary} \textbf{\textcolor{red!70!yellow}{Analysis}}


Consider the general square root function.

\[  A \sqrt{B \,x + C} + D  \]


\textbf{Domain:} The domain is all real numbers that make the inside nonnegative


\[ B \,x + C >= 0 \]


\textbf{Zeros:} Square root functions can have a single zero or no zeros.


\[  A \sqrt{B \,x + C} + D  = 0 \]

Since square roots of real numbers are always positive (or $0$), this equation has a solution only if $A$ and $D$ have opposite signs.


\textbf{Continuity:} Square root functions are contuous functions.  They are continuous on their domain.



\textbf{End-Behavior:} Square root functions have end-behavior only in the direction that makes the inside positive. In that direction, the value of the function approaches $-\infty$ or $\infty$, which is dictated by the sign of the leading coeffcient.



\textbf{Behavior:} Square root functions are either increasing functions or they are decreasing functions.  

The Chain Rule will give us this information from the behavior of its component functions.


\[
(L_{out} \circ SQ \circ L_{in})(x) = A \sqrt{B \,x + C} + D
\]


$SQ$ is the core square root function and is an increasing function.  $L_{out}$ and $L_{in}$ are linear functions and their behavior is given by the sign of their leading coefficients.

Then apply the Chain Rule.



\textbf{Global Maximum and Minimum:}  Square root functions have either a global minimum or a global maximum, but not both. This global extreme value occurs at the real number that makes the inside of the radical equal to $0$.


\textbf{Local Maximums and Minimums:} Same as global.


\textbf{Range:} The range is either of the form $[m, \infty)$ or $(-\infty, M]$.

\end{summary}



























\subsection*{Odd Roots}

This course is the study of the real numbers. As a result we don't take even roots of negative numbers, however we do have odd roots of negative numbers

Odd roots

\[   Odd(x) = \sqrt[2n+1]{x} = x^{\tfrac{1}{2n+1}}          \]

look a lot like the cube root


\[   Odd(x) = \sqrt[3]{x} = x^{\tfrac{1}{3}}          \]



Their domains include all real numbers: $(-\infty, \infty)$.  Odd root functions increase very slowly over this domain, becoming unbounded.







The graph of $y = CubeRoot(x) = \sqrt[3]{x}$


\begin{image}
\begin{tikzpicture} 
  \begin{axis}[
            domain=-10:10, ymax=10, xmax=10, ymin=-10, xmin=-10,
            axis lines =center, xlabel=$x$, ylabel=$y$, 
            ytick={-10,-8,-6,-4,-2,2,4,6,8,10},
            xtick={-10,-8,-6,-4,-2,2,4,6,8,10},
            ticklabel style={font=\scriptsize},
            every axis y label/.style={at=(current axis.above origin),anchor=south},
            every axis x label/.style={at=(current axis.right of origin),anchor=west},
            axis on top
          ]
          
            %\addplot [line width=2, penColor, smooth, domain=(-9:0),<->] {(x-1)/((x+3)*(x-4))};
            \addplot [line width=2, penColor, smooth, samples=200,domain=(0:9),->] {x^0.33333};
            \addplot [line width=2, penColor, smooth, samples=200,domain=(-9:0),<-] {-(-x)^0.33333};
   
      %\addplot[color=penColor,fill=penColor,only marks,mark=*] coordinates{(0,0)};

           

  \end{axis}
\end{tikzpicture}
\end{image}



The cube root function has a vertical tangent line where the inside of the radical equals $0$.






































All odd root/radical functions can be viewed as the Core Cube Root function composed with linear function. One linear function affects the inside of the radical sign and the other affects the outside.


\begin{idea} \textbf{\textcolor{blue!55!black}{Cube Root Functions}} 


Let $CR(x) = \sqrt[3]{k}$ an increasing function.

Let $L_{in}(t) = B \, t + C$ be a linear function.

Let $L_{out}(y) = A \, y + D$ be a linear function.


In general, cube root (radical) functions are constructed by composing these component functions.


\[
(L_{out} \circ CR \circ L_{in})(x) = A \sqrt[3]{B \,x + C} + D
\]


\end{idea}

















\begin{example} Cube Root Function


The graph of $y = g(t) = -4 \sqrt[3]{2t+3}+1$


\begin{image}
\begin{tikzpicture} 
  \begin{axis}[
            domain=-10:10, ymax=10, xmax=10, ymin=-10, xmin=-10,
            axis lines =center, xlabel=$t$, ylabel=$y$, 
            ytick={-10,-8,-6,-4,-2,2,4,6,8,10},
            xtick={-10,-8,-6,-4,-2,2,4,6,8,10},
            ticklabel style={font=\scriptsize},
            every axis y label/.style={at=(current axis.above origin),anchor=south},
            every axis x label/.style={at=(current axis.right of origin),anchor=west},
            axis on top
          ]
          
            \addplot [line width=2, penColor, smooth, samples=200,domain=(-1.5:8),->] {-4*(2*x+3)^0.33333+1};

            \addplot [line width=2, penColor, smooth, samples=200,domain=(-6:-1.5),<-] {4*(-2*x-3)^0.33333+1};

            %\addplot[color=penColor,fill=penColor,only marks,mark=*] coordinates{(-3,0)};

           

  \end{axis}
\end{tikzpicture}
\end{image}


\textbf{Domain:}  The domain of all cube root functions is $(-\infty, \infty)$.


\textbf{Zeros:} Cube root functions always have one zero.


\begin{align*}
-4 \sqrt[3]{2t+3}+1      & = 0 \\
\sqrt[3]{2t+3}           & = \frac{1}{4} \\
2t+3                     & = \left( \frac{1}{4} \right) = \frac{1}{64}  \
t          & = \frac{1}{2} \left( \frac{1}{64} - 3 \right)
          & = -\frac{191}{128}
\end{align*}


Just Checking...$-\frac{191}{128} \approx 1.49$, which agrees with the graph.


\textbf{Continuity:} Cube root functions are continuous.



\textbf{End-Behavior:} Cube root functions are unbounded in either direction.  The behavior will tell us which way.




\textbf{Behavior:} 


$g$ is the composition of

\begin{itemize}
\item $CR(x) = \sqrt[3]{x}$, an increasing function
\item $L_{in}(t) = 2t+3$, an increasing linear function
\item $L_{out}(y) = -4y + 1$, a decreasing linear function
\end{itemize}


\[ g(t) = -4 \sqrt[3]{2t+3}+1 =  (L_{out} \circ g \circ L_{in})(t) \]

\[ dec \circ inc \circ inc = dec  \]


$g$ is a decreasing cube root function, which agrees with the graph.

This also gives us the end-behavior.


\[ \lim\limits_{t \to -\infty} g(t) = \infty \, \text{ and } \,  \lim\limits_{t \to \infty} g(t) = -\infty \]



\textbf{Global Maximum and Minimum:}  Cube root functions do not have global a maximum or minimum. 


\textbf{Local Maximums and Minimums:}  Cube root functions do not have local maximums or minimums. 


\textbf{Range:} The range of all cube root functions is $(-\infty, \infty)$. 




\end{example}
















\begin{summary} \textbf{\textcolor{red!70!yellow}{Analysis}}


Consider the general cube root function.

\[  A \sqrt[3]{B \,x + C} + D  \]


\textbf{Domain:} The domain is all real numbers.


\textbf{Zeros:} Cube root functions always have a single.


\[  A \sqrt[3]{B \,x + C} + D  = 0 \]




\textbf{Continuity:} Cube root functions are contuous functions.   



\textbf{End-Behavior:} Cube root functions have opposite end-behavior on the two sides.  This is dictated by the behavior of the function.



\textbf{Behavior:} Cube root functions are either increasing functions or they are decreasing functions.  

The Chain Rule will give us this information from the behavior of its component functions.


\[
(L_{out} \circ CR \circ L_{in})(x) = A \sqrt[3]{B \,x + C} + D
\]


$CR$ is the core cube root function and is an increasing function.  $L_{out}$ and $L_{in}$ are linear functions and their behavior is given by the sign of their leading coefficients.

Then apply the Chain Rule.



\textbf{Global Maximum and Minimum:}  Cube root functions do not have a global maximum or minimum.


\textbf{Local Maximums and Minimums:} Cube root functions do not have a local maximums or minimums.


\textbf{Range:} The range is always $(-\infty, \infty)$.

\end{summary}




































\begin{center}
\textbf{\textcolor{green!50!black}{ooooo-=-=-=-ooOoo-=-=-=-ooooo}} \\

more examples can be found by following this link\\ \link[More Examples of Analysis]{https://ximera.osu.edu/csccmathematics/precalculus/precalculus/radicalFunctions/examples/exampleList}

\end{center}









\end{document}
